%! Exponential Functions — Unit 4, Lesson 4.3 — Warm-Up
\documentclass[10pt]{article}
\usepackage{precalculus-article}
\usepackage{precalculus-boxes}
\newcommand{\TallMath}[1]{$\displaystyle #1\rule[-1.4em]{0pt}{3.2em}$}

\begin{document}

\pageheader{Unit 4, Lesson 4.3}{Warm-Up}
\namedateperiod

\vspace{0.14in}

\begin{enumerate}[itemsep=16pt]

\item A table of values is shown below.

      \renewcommand{\arraystretch}{1.4}
      \begin{tabular}{|c|c|c|c|c|}
        \hline
        $x$    & 0 & 1  & 2  & 3  \\
        \hline
        $f(x)$ & 8 & 12 & 18 & 27 \\
        \hline
      \end{tabular}

      \vspace{0.08in}
      Compute the ratio of each successive output: $\dfrac{12}{8} = $ \blank{1.5cm},
      $\dfrac{18}{12} = $ \blank{1.5cm},
      $\dfrac{27}{18} = $ \blank{1.5cm}

      \vspace{0.06in}
      Are the ratios constant? \blank{2cm} \quad
      Is the function \textbf{linear} or \textbf{exponential}? Circle one.

\item Write the exponential function $f(x) = ab^x$ that passes through $(0,\,5)$ with
      common ratio $b = 3$.

      \vspace{0.06in}
      $f(x) = $ \blank{6cm}

\item Write the linear function $g(x) = mx + c$ that passes through $(2,\,11)$ and $(5,\,23)$.

      \vspace{0.06in}
      Slope: $m = \dfrac{23 - 11}{5 - 2} = $ \blank{2cm}

      \vspace{0.06in}
      $g(x) = $ \blank{6cm}

\end{enumerate}

\end{document}
